\documentclass[12pt,a4paper]{article}

\usepackage[utf8]{inputenc}
\usepackage[T1]{fontenc}
\usepackage{lmodern}
\usepackage[margin=1in]{geometry}
\usepackage{hyperref}
\usepackage{setspace}
\usepackage[indonesian]{babel}

\hypersetup{
    colorlinks=true,
    linkcolor=blue,
    citecolor=blue,
    urlcolor=blue
}

\title{\textbf{Ketika Kesepakatan Kehilangan Landasannya}}
\author{Teguh Wicaksono\thanks{Penulis adalah dosen Fakultas Ekonomi dan Bisnis, Universitas Islam Internasional Indonesia (UIII).}}
\date{}

\begin{document}

\maketitle
\thispagestyle{empty}

\onehalfspacing

Dani Rodrik pernah menulis bahwa negara berkembang menghadapi ``trilemma'': kita tidak bisa secara bersamaan mempertahankan kedaulatan kebijakan domestik, integrasi ekonomi mendalam, dan ruang fiskal yang memadai. Kita harus memilih dua dari tiga.

Dalam sepuluh hari terakhir, Indonesia seolah menguji batas trilemma itu secara dramatis.

Mari kita rekonstruksi kronologinya. Sejak Juli 2025, ketika Trump mengancam tarif 32 persen melalui IEEPA, tim negosiasi Indonesia bekerja di bawah tekanan luar biasa. Tarif diturunkan menjadi 19 persen per Agustus. Namun Washington menginginkan lebih. Negosiasi berlanjut berbulan-bulan hingga puncaknya pada 19 Februari 2026, ketika Presiden Prabowo dan Presiden Trump menandatangani \textit{Agreement of Reciprocal Trade} di Washington.

Untuk mendapatkan 19 persen dan pembebasan 1.819 pos tarif, Indonesia---seperti dilaporkan Jakarta Post---``membalikkan sejumlah kebijakan yang dipertahankan bertahun-tahun.''

Daftarnya panjang. TKDN---persyaratan kandungan lokal yang selama ini menjadi instrumen industrialisasi---dikecualikan untuk perusahaan AS. Sertifikasi halal dilonggarkan untuk produk asal AS. Yang paling mengkhawatirkan dari perspektif kedaulatan digital: Indonesia mengakui AS sebagai yurisdiksi dengan ``perlindungan data memadai,'' mengizinkan transfer data pribadi warga Indonesia ke luar negeri---meskipun AS sendiri tidak memiliki undang-undang perlindungan data federal setara GDPR Eropa. Restriksi ekspor mineral kritis dilonggarkan, menggerus pilar kebijakan hilirisasi. Pasar pertanian dibuka lebar: 50.000 ton daging sapi AS per tahun, satu juta ton gandum, penghapusan hambatan untuk kedelai dan jagung. Izin Freeport-McMoRan di Grasberg diperpanjang 20 tahun hingga 2061. Dan ada Pasal 5.1 yang nyaris luput dari perhatian: Indonesia harus ``mencerminkan pembatasan perdagangan AS terhadap negara ketiga terkait keamanan nasional Amerika''---berpotensi memaksa Indonesia mengikuti sanksi Washington.

Pertanyaannya: apakah semua ini sepadan?

Dalam teori negosiasi, kita mengenal konsep BATNA---\textit{Best Alternative to a Negotiated Agreement}. Kekuatan posisi tawar ditentukan oleh seberapa baik alternatif Anda jika negosiasi gagal. Ekspor Indonesia ke AS hanya 9,9 persen dari total ekspor kita. Surplus perdagangan dengan AS kurang dari 2 persen PDB. Bandingkan dengan Vietnam: 29,4 persen ekspor ke AS, surplus 21,8 persen PDB. Vietnam punya alasan lebih kuat untuk memberikan konsesi besar. Indonesia? Kurang jelas.

Namun hal yang paling dramatis terjadi kemudian. Dua puluh empat jam setelah penandatanganan, Mahkamah Agung AS dalam putusan 6--3 menyatakan bahwa IEEPA tidak memberikan wewenang kepada presiden untuk mengenakan tarif. Landasan hukum seluruh tarif resiprokal---dan seluruh kesepakatan bilateral yang dinegosiasikan di bawahnya---runtuh.

Bayangkan Anda membeli rumah. Anda sudah menandatangani kontrak, membayar uang muka, bahkan mulai merenovasi. Lalu pengadilan menyatakan bahwa penjual tidak pernah memiliki hak jual atas properti itu.

Inilah posisi Indonesia sekarang.

Pemerintah Trump segera mengenakan tarif pengganti 10 persen berdasarkan Pasal 122 UU Perdagangan 1974---landasan hukum berbeda, dengan keterbatasan sangat berbeda pula. Pasal 122 tidak bisa membedakan antarnegara: semua menghadapi tarif sama. Dan ada batas waktu: 150 hari tanpa perpanjangan Kongres.

Mari kita lihat datanya. Dengan menghitung tarif efektif tertimbang berdasarkan \$28,05 miliar ekspor Indonesia pada 2024, saya menemukan: tarif efektif pasca-putusan (15,3 persen) justru \textit{lebih rendah} dari tarif efektif di bawah perjanjian bilateral (19,7 persen). Penurunan 4,4 poin persentase ini terjadi karena tarif Pasal 122 yang 10 persen lebih rendah dari 19 persen resiprokal.

Namun ada ironi pahit. Sektor-sektor yang pembebasan tarifnya diperjuangkan melalui negosiasi panjang---minyak sawit, karet, kopi, kakao---kehilangan perlindungan. Minyak sawit naik dari 2 persen menjadi 12 persen. Karet dari 3 menjadi 13 persen. Sementara tekstil dan alas kaki justru diuntungkan: tarif turun dari 29 persen menjadi 22 persen tanpa perlu negosiasi apa pun.

Yang menarik adalah ini: Indonesia sekarang menghadapi tarif yang \textit{sama} dengan Vietnam, Bangladesh, dan Uni Eropa. Di bawah perjanjian bilateral, Indonesia justru membayar 19 persen---lebih tinggi dari tarif 10 persen yang dikenakan pada negara-negara yang \textit{belum} menyepakati \textit{deal}. Secara paradoks, putusan Mahkamah Agung menyamakan kedudukan Indonesia.

Tapi paradoks ini memunculkan pertanyaan lebih mendasar. Indonesia telah memberikan konsesi besar---TKDN, halal, data pribadi, mineral, Freeport, pembelian \$33 miliar---sebagai imbalan atas akses pasar yang terstruktur. Akses itu kini tidak dapat diberikan sebagaimana dijanjikan. Prinsip resiprositas, fondasi hukum perdagangan internasional, menyarankan bahwa kewajiban Indonesia perlu ditinjau ulang.

Persoalannya tidak sesederhana ``batalkan semuanya.'' Banyak konsesi berupa perubahan kebijakan struktural yang, bila dibatalkan, merusak kepercayaan investor. Moody's sudah menurunkan prospek kredit Indonesia. MSCI memperingatkan kemungkinan penurunan status pasar ekuitas. Membatalkan komitmen secara sepihak bukan tanpa risiko. Namun mempertahankan seluruh konsesi tanpa imbalan setara juga bukan pilihan yang masuk akal.

Maka yang diperlukan adalah kehati-hatian yang cerdas. \textit{Pertama}, tunda ratifikasi DPR hingga kerangka hukum baru di Washington terbentuk---perjanjian memerlukan proses domestik 90 hari, gunakan waktu itu. \textit{Kedua}, ajukan klaim pengembalian bea masuk yang dibayar di bawah IEEPA---potensinya lebih dari \$500 juta. \textit{Ketiga}, manfaatkan jendela 150 hari di mana semua pesaing menghadapi tarif sama untuk merebut kembali pangsa pasar.

Dan pelajaran yang lebih besar: ketergantungan pada satu pasar, betapapun besarnya, menciptakan kerentanan eksistensial. RCEP, negosiasi CEPA dengan Uni Eropa, pendalaman hubungan dengan Timur Tengah dan Afrika---bukan lagi opsi, melainkan keharusan.

Indonesia tidak perlu euforia karena tarif agregat turun, tidak perlu pula panik karena pembebasan komoditas hilang. Yang diperlukan adalah kejelasan berpikir: apa yang sudah kita berikan, apa yang masih bisa kita pertahankan, dan apa yang harus kita minta kembali.

Dalam negosiasi, seperti kata Thomas Schelling, kekuatan terbesar kadang justru muncul dari kemampuan untuk \textit{tidak} berkomitmen terlalu cepat. Mungkin itulah pelajaran paling mahal dari 19 Februari.

\end{document}
